\documentclass{report}
\usepackage{graphicx} % Required for inserting images

\usepackage{fontspec}
\setmainfont{FreeSerif}

\usepackage{polyglossia}
\setdefaultlanguage{greek}
\setotherlanguages{english}
\usepackage[normalem]{ulem}

\usepackage{caption}
\usepackage[useregional]{datetime2}
\usepackage{amsmath}

\usepackage{url}
\usepackage{hyperref}
\usepackage{biblatex}
\bibliography{citation}
%\addbibresource{citation}
\usepackage[dvipsnames]{xcolor}

\hypersetup{
	colorlinks=true,
	linkcolor=black,
	filecolor=magenta,      
	urlcolor=cyan,
	citecolor = violet
}

\urlstyle{same}

\DTMlangsetup{showdayofmonth=false}

\title{Σύγκριση Αναμενόμενης Επιτάχυνσης για τιμές $V_{s30}$ και Κλίσης-Γεωλογίας\\
}
\author{Διαμαντής Χατζηαργυρίου}
\date{\Today}

\begin{document}

\maketitle

\section{Εισαγωγή}
Σε συνέχεια της προηγούμενης προκαταρκτικής ανάλυσης, γίνεται σύγκριση των αποτελεσμάτων μεταξύ μεθόδων.

Εν προκειμένω, γίνεται πιθανότική ανάλυση της σεισμικής επικινδυνότητας για διαφορετικές μεθόδους, χρησιμοποιώντας από την μία παραμέτρους της κλίσης και γεωλογίας, και από την άλλη την $V_{s30}$.

\section{Ανάλυση}
%\subsection{Μοντέλα}
Δύο σκέλη αναλύσεων πραγματοποιούνται και συγκρίνονται. Μία για ανάλυση με $V_{s30}$, όπως προκύπτει από παρεμβολή σε γεωτρήσεις. Και μία δεύτερη, που χρησιμοποιεί την γεωλογία της θέσης και την κλίση της επιφάνειας σε $mm/mm$, όπως αυτή προκύπτει από το εργαλείο ESRM20\_sitemodel.

Τα σημεία που παρουσιάζονται προκύπτουν από το API του OSM κατά μήκος της πρώτης γραμμής και πρώτης τροχιάς του μετρό Θεσσαλονίκης, και πυκνώθηκαν γραμμικώς από επιπλέον αλγόριθμο. Έτσι, βελτιώνεται η ευκρίνεια των αποτελεσμάτων.

Η επιρροή της πύκνωσης φαίνεται ιδιαιτέρως στις ανατεθιμένες τιμές της $V_{s30}$, όπου για σημεία κοντινά το ένα στο άλλο, η τιμή παραμένει σταθερή. Η αιτία να είναι η μέθοδος που το ESRM αναθέτει τις παραμέτρους για ευρύτερες περιοχές.

Η ανάθεση τιμών $V_{s30}$ στις θέσεις μελέτης έγινε με παρεμβολή των αντίστοιχων τιμών σε εγγείς γεωτρήσεις, αξιοποιώντας γεωστατιστικό αλγόριθμο Kriging με εύρος 330.

Ο συγκεκριμένος αλγόριθμος δεν έχει μελετηθεί σε βάθος στην συγκεκριμένη εργασία, και πιθανώς να χρειάζεται περαιτέρω βαθμονόμηση για αξιόπιστα αποτελέσματα.

Οι δύο αναλύσεις πραγματοποιούνται για πανομοιότυπες θέσεις παραμέτρους, εξαιρουμένες τις παραμέτρους πεδίου και τα λογικά δέντρα.

\section{Αποτελέσματα}
Στο Σχήμα \ref{fig:anal} αναγράφονται τα αποτελέσματα της ανάλυσης.

Σε αυτά διακρίνεται ότι η ανάλυση για τις συγκεκριμένες $V_{s30}$ σε σχέση με αυτήν για κλίσεις δίνει μεγαλύτερη μέση PGA. Αντίστοιχα, η μέση PGA για δεδομένα ESRM είναι πανομοιότυπη με αυτήν που είχε υπολογιστεί στην πρώτη ανάλυση.

Μία απαραίτητη παρατήρηση είναι πως η γραφική συσχέτιση που υπάρχει μεταξύ PGA και $V_{s30}$ δεν ανταποκρίνεται στην πραγματικότητα, καθώς μόνο η μία μέθοδος χρησιμοποιεί την συγκεκριμένη παράμετρο. Μία νέα ανάλυση θα έπρεπε να παραχθεί για την συγκεκριμένη συσχέτιση.

\begin{figure}[ht]
	\centering	
	\includegraphics[width=\linewidth]{../data/Plot3.png}
 	\caption{Αποτελέσματα Ανάλυσης}
	\label{fig:anal}
\end{figure}

Εξαιρετικό ενδιαφέρον όμως, παρουσιάζουν οι καμπύλες επικινδυνότητας. Ενώ η μέση τιμή της PGA είναι σαφώς μεγαλύτερη για όλα τα σημεία από την ανάλυση για $V_{s30}$, οι καμπύλες επικινδυνότητας στις δύο αναλύσεις μένουν σχεδόν πανομοιότυπες.

Για παράδειγμα, απεικονίζονται οι καμπύλες για τον Ν.Σ.Σ., όπου η διαφορά σε μέση PGA είναι πολύ υψηλή. Παρ' όλα αυτά, η απόκλιση στις καμπύλες είναι ιδιαιτέρως χαμηλή, παρουσιάζοντας εμφανείς διαφορές σε εξαιρετικά υψηλές περιόδους επαναφοράς γεγονότος (μεγαλύτερες των 10,000 ετών).

Φυσικά, το διάγραμμα είναι λογαριθμικό, αλλά ακόμα κι έτσι η διαφορά στις μεθόδους δεν φαίνεται να είναι τόσο επιδραστική στα αποτελέσματα.

\end{document}
